\chapter{4x10\textsuperscript{9}D.C.}\label{4x10^9-4}

\hfill\break

Attraversai qualche metro di terra pressata su cui l'asfalto segnato
dalle intemperie era attaccato a chiazze scabre, arrivai a un argine e
scivolai di sotto facendo un gran baccano, con le valigie rigide stipate
di vestiti modesti e appunti scritti a mano e file digitali e farmaci
marziani. Andai a finire in una fossa di drenaggio, immerso fino alle
cosce in un'acqua verde come le foglie di papaya e calda come le notti
tropicali. L'acqua rifletteva la Luna deturpata e puzzava di letame.

Nascosi i bagagli in un luogo asciutto a metà dell'argine e mi trascinai
fino in cima, sdraiato in un angolo che mi nascondeva ma mi consentiva
di vedere la strada, la clinica quadrata di cemento di Ibu Ina e l'auto
nera parcheggiata davanti.

Gli uomini della macchina fecero irruzione dalla porta sul retro. Mentre
si spostavano nell'edificio, accesero altre luci, creando quadrati
gialli con le finestre dalle tende tirate, ma non avevo modo di sapere
cosa stessero facendo. Probabilmente una perquisizione. Cercai di
stimare da quanto erano dentro, ma mi sembrava di avere perso la
capacità di calcolare il tempo o perfino di leggere i numeri
sull'orologio. Le cifre brillavano come lucciole inquiete ma non si
fermavano il tempo sufficiente per assumere un senso.

Uno degli uomini uscì dalla porta d'ingresso, si avvicinò all'auto e
avviò il motore. Il secondo uomo comparve dopo qualche secondo e
s'infilò dalla parte del passeggero. La macchina scura come la notte
sbandò mentre si girava sulla strada nella mia direzione e i fari
inondarono di luce la berma. Abbassai la testa e rimasi immobile fino a
quando il rumore del motore svanì.

Poi mi chiesi cosa avrei potuto fare, a quel punto. Era difficile
trovare una risposta a quella domanda, perché mi sentivo stanco, di una
stanchezza improvvisa, enorme; ero troppo debole per alzarmi. Volevo
tornare alla clinica, trovare un telefono e avvisare Ibu Ina riguardo
agli uomini in macchina. Forse, però, l'avrebbe fatto En. Ci speravo.
Perché non ce l'avrei fatta ad arrivare alla clinica. Volevo muovere le
gambe ma mi tremavano. Non era solo affaticamento, sembrava fossero
paralizzate.

E quando guardai di nuovo verso la clinica, c'era del fumo che
fuoriusciva a spirale dalla canna fumaria del tetto e la luce dietro le
tende era di un giallo scintillante. Fuoco.

I due uomini avevano incendiato la clinica di Ibu Ina, e non c'era nulla
che potessi fare se non chiudere gli occhi e sperare di non morire lì
prima che qualcuno mi trovasse.

ooo

Mi svegliarono la puzza di fumo e il suono di un pianto.

Non era ancora giorno. Mi accorsi di riuscire a muovermi, almeno un po',
con grande sforzo e dolore, e mi sembrò di riuscire a pensare con una
certa lucidità. E così mi rialzai a fatica sul pendio, centimetro dopo
centimetro.

C'erano macchine e persone ovunque, tra l'argine e la clinica, fari e
torce che tracciavano archi spasmodici nel cielo. La clinica era una
rovina fumante. Le pareti di cemento erano ancora in piedi ma il tetto
era crollato e l'edificio era stato eviscerato dal fuoco. Riuscii ad
alzarmi. Camminai in direzione del pianto.

Proveniva da Ibu Ina. Era seduta su un'isola di asfalto con le ginocchia
strette a sé, circondata da un gruppo di donne che mentre mi avvicinavo
mi rivolsero occhiate cupe e sospettose. Ma quando Ina mi vide, balzò in
piedi, asciugandosi gli occhi con la manica. «Tyler Dupree!» Corse verso
di me. «Pensavo fossi morto! Bruciato insieme a tutto il resto!»

Mi afferrò, mi abbracciò, mi sostenne - le gambe mi erano ridiventate
gommose. «La clinica» riuscii a dire. «Tutto il tuo lavoro. Mi dispiace
tanto, Ina\ldots»

«No, la clinica è solo un edificio. L'attrezzatura medica si può
sostituire. Tu, invece, sei unico. En ci ha raccontato tutto, che lo hai
mandato via quando sono arrivati i piromani. Gli hai salvato la vita,
Tyler!» Indietreggiò. «Tyler? Stai bene?»

Non stavo bene. Guardai il cielo alle sue spalle. Era quasi l'alba.
L'antico Sole stava sorgendo. Contro il cielo indaco si delineava il
profilo del monte Merapi. «Sono solo stanco» dissi, e chiusi gli occhi.
Mi sentii cedere le gambe e udii Ina chiedere aiuto, e poi dormii ancora
un po' - per giorni, mi dissero in seguito.

ooo

Per ovvie ragioni, non potevo rimanere nel villaggio.

Ina voleva curarmi per via dell'ultima crisi dovuta al farmaco ed era
convinta che il villaggio dovesse proteggermi. Dopo tutto, avevo salvato
la vita di En (o così insisteva a dire) ed En non era solo suo nipote,
ma di fatto era imparentato con tutto il resto del villaggio, in un modo
o nell'altro. Ero un eroe, ma anche una calamita per gli uomini malvagi;
e se non ci fosse stata Ina a difendermi, a quell'ora suppongo che il
\emph{kepala desa} mi avrebbe caricato sul primo autobus per Padang e
spedito al diavolo. Perciò mi portarono assieme ai miei bagagli in una
casa disabitata (i proprietari si erano trasferiti per \emph{rantau}
mesi prima) per il tempo necessario a trovare un'altra sistemazione.

I Minangkabau della zona occidentale di Sumatra sapevano come chinare la
testa e schivare i colpi dell'oppressione. Erano sopravvissuti alla
venuta dell'Islam nel sedicesimo secolo, alla guerra dei Paderi, al
colonialismo olandese, al Nuovo Ordine di Suharto, alla Restaurazione
Nagari e, dopo lo Spin, alla nuova Reformasi e alla loro criminosa
polizia nazionale. Ina mi aveva raccontato alcune di quelle storie, sia
alla clinica che in seguito, mentre ero disteso in una stanzetta di una
casa di legno sotto le pale enormi e lente di un ventilatore elettrico.
La forza dei Minangkabau, diceva, era la loro flessibilità, la profonda
consapevolezza che il resto del mondo non era la loro casa e non lo
sarebbe mai stata. (Citò un loro proverbio: ``\emph{In campi diversi,
	cavallette diverse; in stagni diversi, pesci diversi}.'') La tradizione
del \emph{rantau} - l'emigrazione di giovani che si spostano in varie
parti del mondo e tornano più ricchi o più saggi - li aveva resi persone
sofisticate. Le semplici case in legno del villaggio, con tetti a forma
di corno di bufalo, erano adornate di antenne aerostatiche e la maggior
parte delle famiglie, disse Ina, riceveva regolarmente lettere o e-mail
dai propri cari sparsi per l'Australia, l'Europa, il Canada, gli Stati
Uniti.

Non c'era da sorprendersi, quindi, che sulle banchine di Padang ci
fossero Minangkabau che lavoravano a ogni livello. L'ex marito di Ina,
Jala, era solo uno dei tanti nel settore import-export che organizzava
spedizioni \emph{rantau} per raggiungere l'Arco o andare oltre. Non era
un caso che le indagini di Diane l'avessero condotta a Jala e quindi a
Ibu Ina e a questo villaggio dell'altopiano. «Jala è un opportunista e
sa essere piuttosto meschino se vuole, ma non è senza scrupoli» disse
Ina. «Diane è stata fortunata a trovarlo, oppure sa giudicare il
carattere delle persone, il che è più probabile. In ogni caso Jala non
nutre alcuna simpatia per la nuova Reformasi, per fortuna di tutti
quelli coinvolti.»

(Aveva divorziato da Jala, mi raccontò, perché aveva preso il brutto
vizio di dormire con donne disdicevoli della città. Spendeva troppi
soldi per le sue ragazze e in due occasioni le aveva portato in casa
malattie veneree curabili ma preoccupanti. Era un cattivo marito, disse
Ina, ma non particolarmente cattivo come uomo. Non avrebbe denunciato
Diane alle autorità a meno che non fosse stato catturato e torturato
fisicamente\ldots{} ma era troppo intelligente per lasciarsi catturare.)

«Gli uomini che hanno bruciato la clinica\ldots»

«Devono aver seguito Diane fino all'hotel di Padang e interrogato
l'autista che ti ha portato qui.»

«Ma perché bruciare l'edificio?»

«Non lo so, ma sospetto che sia stato un tentativo di spaventarti e
farti uscire allo scoperto. E un avvertimento per chiunque voglia
aiutarti.»

«Se hanno trovato la clinica, allora sanno come ti chiami.»

«Ma non verranno nel villaggio alla luce del giorno, ad armi spianate.
Le cose non sono precipitate fino a quel punto. Mi aspetto che stiano lì
a osservare il porto e a sperare che facciamo qualcosa di stupido.»

«Comunque, se il tuo nome è in una lista, se cerchi di aprire un'altra
clinica\ldots»

«Non ne ho mai avuta l'intenzione.»

«No?»

«No. Sei riuscito a convincermi che il \emph{rantau gadang} per un
medico può essere una buona opportunità. Sempre che non ti preoccupi la
concorrenza\ldots»

«Non capisco.»

«Intendo dire che c'è una soluzione semplice per tutti i nostri
problemi, e ce n'è una, in particolare, che sto valutando da un po'.
Tutto il villaggio l'ha presa in considerazione, in un modo o
nell'altro. Molti se ne sono già andati. La nostra non è una città
grande e prospera come Belubus o Batusangkar. La terra qui non è
particolarmente ricca e ogni anno perdiamo persone che se ne vanno in
città o si uniscono ai clan di altre città o al \emph{rantau gadang}, e
perché no? C'è spazio per tutti nel nuovo mondo.»

«Vuoi emigrare?»

«Io, Jala, mia sorella e sua sorella e i miei nipoti e cugini - più di
trenta, in tutto. Jala ha diversi figli illegittimi che sarebbero felici
di assumere il controllo della sua attività una volta che lui è
dall'altra parte. Quindi, vedi?» Sorrise. «Non devi ringraziarci di
nulla. Non siamo i tuoi benefattori. Solo compagni di viaggio.»

Le chiesi diverse volte se Diane fosse al sicuro. Al sicuro per quanto
avrebbe potuto tenerla al sicuro Jala, rispose Ina. Jala l'aveva
sistemata in uno spazio abitabile sopra un ufficio della dogana dove
sarebbe stata abbastanza comoda e nascosta finché non fossero stati
fatti i preparativi finali. «La parte difficile sarà portarti al porto
senza che ti vedano. La polizia sospetta che tu sia nell'altopiano e
controllerà le strade per individuare gli stranieri, soprattutto quelli
malati, perché l'autista che ti ha portato alla clinica gli avrà detto
che non stai bene.»

«Non sto più male» ribattei.

L'ultima crisi era iniziata quando ero fuori della clinica in fiamme ed
era passata mentre ero incosciente. Ibu Ina disse che era stata una
transizione difficile, che dopo il trasferimento in quella piccola
stanza di quella casa vuota mi ero lamentato tanto che i vicini avevano
protestato, che aveva avuto bisogno di suo cugino Adek per tenermi fermo
durante le convulsioni peggiori - ecco perché avevo le braccia e le
spalle piene di lividi, non me n'ero accorto? Ma non mi ricordavo
niente. Ero solo consapevole di sentirmi più forte man mano che
passavano i giorni; la temperatura era nella norma; riuscivo a camminare
senza tremare.

«E gli altri effetti del farmaco?» chiese Ina. «Ti senti diverso?»

Era una domanda interessante. Risposi con sincerità: «Non lo so. Non
ancora, comunque.»

«Bene. Per il momento non ha molta importanza. Come dicevo, la
difficoltà sarà farti uscire dall'altopiano e tornare a Padang. Per
fortuna, credo che riusciremo a organizzare il tutto.»

«Quando ce ne andremo?»

«Fra tre o quattro giorni,» disse Ina, «nel frattempo riposati.»

ooo

Ina fu occupata per gran parte di quei tre giorni. La vidi pochissimo.
Le giornate erano calde e soleggiate, ma la brezza attraversava la casa
in legno portando folate confortanti, e io passavo il tempo a fare
esercizio fisico con prudenza, a scrivere e a leggere - c'erano dei
tascabili in inglese su uno scaffale di vimini nella stanza da letto,
tra cui una biografia in voga di Jason Lawton dal titolo ``\emph{Una
	vita per le stelle}''. (Cercai il mio nome nell'indice e lo trovai:
``Dupree, Tyler'', con cinque riferimenti di pagina. Ma non me la
sentivo di leggere il libro. I romanzi del lordotico Somerset Maugham
erano più allettanti.)

En veniva a trovarmi regolarmente per vedere se stessi bene e mi portava
panini e acqua imbottigliata del \emph{warung} di suo zio. Sfoggiava un
atteggiamento da padrone e si premurava sempre di chiedere notizie sulla
mia salute. Disse di essere «orgoglioso di fare \emph{rantau}» con me.

«Anche tu, En? Andrai nel nuovo mondo?»

Annuì energicamente. «Anche mio padre, mia madre, mio zio» e un'altra
decina di parenti stretti di cui definì i gradi di parentela con parole
minang. Gli brillavano gli occhi. «Forse là m'insegnerai la medicina.»

Forse avrei dovuto. Attraversare l'Arco avrebbe reso quasi impossibile
un'istruzione tradizionale. Forse non era la cosa migliore per En, e mi
chiesi se i suoi genitori avessero riflettuto a sufficienza sulla loro
decisione.

Ma non era affar mio, ed En era visibilmente euforico per il viaggio.
Quando ne parlava quasi non riusciva a controllare la voce. Mi piaceva
l'espressione entusiasta e impaziente sul suo viso. En apparteneva a una
generazione capace di guardare al futuro con speranza più che timore.
Nessuno di quelli della mia età aveva mai sorriso al futuro in quel
modo. Era uno sguardo buono, profondamente umano, che mi rendeva felice
e mi rendeva triste.

Ina tornò la sera prima della partenza. Portò la cena e un piano.

«Il cognato del figlio di mio cugino guida le ambulanze per l'ospedale
di Batusangkar. Ne può prendere in prestito una dal parco macchine per
portarti a Padang. Ci precederanno almeno due auto con telefoni
wireless, per cui se ci fosse un blocco stradale ci avviserebbero.»

«Non ho bisogno di un'ambulanza» protestai.

«L'ambulanza è una finzione. Tu starai nel retro, nascosto, io sarò
vestita da medico e uno del villaggio - En mi sta supplicando di dargli
quella parte - interpreterà il malato. Capisci? Se la polizia guarda nel
retro dell'ambulanza vedrà me e un bambino malato, e se io dico ``SDC''
gli agenti saranno restii ad approfondire l'ispezione. E così il medico
americano esageratamente alto passerà senza che nessuno se ne accorga.»

«Pensi che funzionerà?»

«Penso che le probabilità di riuscita siano molto buone.»

«Ma se ti trovano con me\ldots»

«Anche nella peggiore delle ipotesi, la polizia non può arrestarmi se
non ho commesso un crimine. Trasportare un occidentale non è un
crimine.»

«Ma trasportare un criminale potrebbe esserlo.»

«Sei un criminale, Pak Tyler?»

«Dipende da come interpreti certi atti del Congresso.»

«Preferisco non interpretarli affatto. Per favore, non preoccupartene.
Ti ho detto che il viaggio è stato rimandato di un giorno?»

«Perché?»

«Per via di un matrimonio. Naturalmente, non sono più quelli di una
volta. Gli \emph{adat} del matrimonio sono stati terribilmente intaccati
da quando c'è lo Spin. Come ogni altra cosa da quando sull'altopiano
sono arrivati i soldi, le strade e i fast food. Non credo che il denaro
sia il male, ma può essere molto corrosivo. I giovani oggigiorno vanno
di fretta. Almeno non abbiamo i matrimoni lampo stile Las Vegas\ldots{}
Esistono ancora nel tuo paese?»

Ammisi di sì.

«Be', pure noi stiamo andando in quella direzione. \emph{Minang hilang,
	tinggal kerbau}. Almeno ci sarà ancora un \emph{palaminan}, un sacco di
riso appiccicoso e musica \emph{saluang}. Te la senti di partecipare?
Almeno per la musica?»

«Ne sarei onorato.»

«Così domani sera cantiamo e la mattina dopo sfidiamo il Congresso
americano. Anche il matrimonio gioca a nostro favore. Tanti spostamenti,
tante auto per strada; non daremo nell'occhio con il nostro piccolo
gruppo \emph{rantau} diretto a Teluk Bayur.»

Mi addormentai tardi e mi svegliai sentendomi meglio, come non succedeva
da tempo, più forte e leggermente più vigile. La brezza mattutina era
calda e satura dell'odore di cibo in cottura, del canto lamentoso dei
galli e del martellare proveniente dal centro della città dove stavano
costruendo un palco all'aperto. Trascorsi la giornata vicino alla
finestra a leggere e guardare il corteo pubblico della sposa e dello
sposo diretti alla casa di lui. Il villaggio di Ina era così piccolo che
il matrimonio lo aveva paralizzato. Per l'occasione anche i
\emph{warung} locali avevano chiuso, malgrado sulla strada principale i
negozi appartenenti alle catene fossero aperti per i turisti. Nel tardo
pomeriggio l'odore di pollo al curry e di latte di cocco aveva ormai
intriso l'aria; En si fece vivo di sfuggita per portarmi un pasto
pronto.

Ibu Ina, con un abito ricamato e un velo di seta, comparve alla porta
poco dopo l'imbrunire e disse: «È finito, intendo il matrimonio vero e
proprio. Ora non resta che cantare e ballare. Vuoi ancora venire,
Tyler?»

Indossai i vestiti migliori che avevo, pantaloni bianchi di cotone e una
camicia dello stesso colore. Il pensiero di essere visto in pubblico mi
rendeva nervoso, ma Ina mi assicurò che non ci sarebbero stati estranei
alla festa di nozze e che tra la folla sarei stato il benvenuto.

Mentre camminavamo insieme sulla strada verso il palco e la musica,
malgrado le rassicurazioni di Ina, sentivo di dare troppo nell'occhio,
non tanto per la mia altezza quanto perché ero rimasto per troppo tempo
al chiuso. Lasciare la casa fu come passare dall'acqua all'aria;
all'improvviso quello che mi circondava era privo di sostanza. Ina mi
distrasse parlandomi dei novelli sposi. Lo sposo, un apprendista
farmacista di Belubus, era un suo giovane cugino. (Ina chiamava
``cugino'' qualsiasi parente che non fosse fratello, sorella, zia o zio;
il sistema di parentela minang utilizzava termini precisi che non
avevano un equivalente nella nostra lingua.) La sposa era una ragazza
del posto con un passato un po' disdicevole. Dopo il matrimonio,
entrambi avrebbero fatto \emph{rantau}. Il nuovo mondo li chiamava.

La musica cominciò al crepuscolo e sarebbe continuata, disse lei, fino
al mattino. Era trasmessa in tutto il villaggio attraverso enormi
altoparlanti montati su palo, ma la fonte era il palco rialzato e il
gruppo che vi stava seduto su stuoie di canna: due uomini strumentisti e
due donne cantanti. Le canzoni, spiegava Ina, parlavano d'amore,
matrimonio, delusione, destino, sesso. Molto sesso, a cui si alludeva
con metafore che Chaucer avrebbe apprezzato. Ci sedemmo su una panchina
lontana dal cuore dei festeggiamenti. Attirai più di uno sguardo curioso
da parte di alcune persone tra la folla - probabilmente venute a
conoscenza della storia della clinica bruciata e dell'americano
fuggitivo - ma Ina fu attenta a non lasciarmi diventare elemento di
disturbo. Rimase in disparte con me, malgrado sorridesse con indulgenza
ai giovani che assaltavano il palco. «Ho superato l'età delle lamentele.
Non c'è più bisogno di ``\emph{arare il mio campo}'', come dice la
canzone, ma\ldots{} mamma mia, che chiasso!»

La sposa e lo sposo con i loro eleganti abiti ricamati erano seduti su
troni finti vicino alla tribuna. Avevo l'impressione che lo sposo, con
quei suoi baffetti sottili a frusta, sembrasse un tipo ambiguo; invece
no, insisteva Ina, era la ragazza, che pareva così innocente con quel
vestito di broccato azzurro e bianco, quella da tenere d'occhio. Bevemmo
del latte di cocco. Sorridemmo. Avvicinandosi la mezzanotte molte donne
del villaggio a poco a poco si allontanarono, lasciando una folla di
uomini, i giovani che facevano ressa attorno al palco e ridevano, gli
adulti seduti ai tavoli intenti a giocare a carte, facce inespressive
come pelle invecchiata.

Avevo mostrato a Ina le pagine che avevo scritto sul mio primo incontro
con Wun Ngo Wen. «Ma il resoconto non sembra del tutto preciso» disse
lei in un momento di tregua dalla musica. «Sembravi fin troppo calmo.»

«Non ero per niente calmo. Cercavo solo di non farmi prendere
dall'imbarazzo.»

«Dopotutto, ti avevano presentato a un uomo di Marte\ldots» Guardò il
cielo, le stelle del dopo Spin nelle loro fragili costellazioni sparse,
offuscate dal bagliore della festa nuziale. «Cosa ti aspettavi?»

«Credo qualcuno meno umano.»

«Ah, quindi era \emph{molto} umano.»

«Sì.»

Wun Ngo Wen era diventato una personalità venerata nell'India rurale, in
Indonesia, nel Sud-Est asiatico. A Padang, disse Ina, nelle case della
gente qualche volta potevi ritrovarti la sua foto in una cornice dorata,
come un santo dipinto ad acquerello o un famoso mullah. «C'era qualcosa
di davvero attraente nel suo atteggiamento. Un modo familiare di
parlare, anche se lo sentivamo solo in traduzione. E quando abbiamo
visto le fotografie del suo pianeta - tutti quei campi coltivati - ci è
sembrato molto più rurale che urbano. Più orientale che occidentale. La
Terra visitata da un ambasciatore di un altro mondo, ed era uno di noi!
Almeno così pareva. E quanto è stato divertente il modo in cui ha
castigato gli americani.»

«L'ultima cosa che Wun intendeva fare era rimproverare chicchessia.»

«Senz'altro la leggenda supera la realtà. Il giorno che te l'hanno
presentato non avresti voluto fargli mille domande?»

«Certo, ma ho immaginato che da quando era arrivato non avesse fatto
altro che rispondere a domande ovvie. Ho pensato che potesse esserne
stufo.»

«Era restio a parlare della sua terra?»

«No, per niente. Amava parlarne. Non gli piaceva che gli facessero
l'interrogatorio, però.»

«Non ho i modi raffinati che hai tu. Sono certa che l'avrei offeso con
le mie domande infinite. Supponiamo, Tyler, che quel primo giorno avessi
potuto chiedergli qualunque cosa: cosa gli avresti chiesto?»

La risposta era facile. Sapevo esattamente quale domanda avevo soffocato
la prima volta che incontrai Wun Ngo Wen. «Gli avrei chiesto dello Spin.
Degli Ipotetici. Se la sua gente aveva appreso qualcosa che noi non
sapevamo ancora.»

«E ne \emph{hai} poi parlato con lui?»

«Sì.»

«E aveva molto da dire?»

«Molto.»

Diedi un'occhiata al palco. Era arrivato un altro gruppo \emph{saluang}.
Uno di loro suonava un \emph{rabab}, uno strumento a corda. Il musicista
batteva l'archetto sulla pancia del \emph{rabab} e sorrideva. Un'altra
lasciva canzone di nozze.

«Forse l'interrogatorio l'ho fatto io a te» disse Ina.

«Mi dispiace. Sono ancora un po' stanco.»

«Allora dovresti andare a casa a dormire. Ordini del dottore. Con un po'
di fortuna domani rivedrai Ibu Diane.»

Abbandonammo i festeggiamenti percorrendo insieme la strada rumorosa. La
musica proseguì fin quasi alle cinque del mattino dopo. Io comunque
dormii profondamente.

ooo

L'autista dell'ambulanza era un uomo magro e taciturno con la divisa
bianca della Mezzaluna Rossa. Si chiamava Nijon, mi strinse la mano con
una deferenza eccessiva e, mentre mi parlava, continuava a guardare Ibu
Ina con i suoi occhi grandi. Gli chiesi se fosse nervoso per il fatto di
dovermi portare a Padang. Ina tradusse la sua risposta: «Dice di aver
fatto cose più pericolose per ragioni meno impellenti, che è contento di
conoscere un amico di Wun Ngo Wen, e dice anche che dovremmo metterci in
viaggio il prima possibile.»

E così salimmo nella parte posteriore dell'ambulanza. Lungo una parete
c'era un armadietto orizzontale in acciaio dove di solito venivano
riposte le attrezzature. Serviva anche da panca. Nijon lo aveva
svuotato; stabilimmo che potevo entrarci piegando le ginocchia
all'altezza dei fianchi e abbassando la testa sotto una spalla.
L'armadietto puzzava di antisettico e lattice ed era comodo quasi quanto
una cassa acustica, ma era lì che mi sarei dovuto nascondere se ci
avessero fermati a un posto di blocco, con Ina seduta sulla panca col
suo camice da ambulatorio ed En steso su una barella che avrebbe fatto
del suo meglio per dare l'impressione di essere affetto da SDC. Alla
calda luce del mattino il piano sembrava più che leggermente ridicolo.

Nijon aveva infilato delle zeppe sotto il coperchio dell'armadietto per
lasciar circolare un po' d'aria all'interno, così forse non sarei
soffocato, ma la prospettiva di passare del tempo in quella cassa di
metallo buia e calda non mi attirava per niente. Fortunatamente - una
volta stabilito che ci stavo - non fui costretto a entrarci, perlomeno
non ancora. Tutte le operazioni di polizia, disse Ina, si erano
concentrate sulla nuova autostrada tra Bukik Tinggi e Padang, e poiché
viaggiavamo in un convoglio sparpagliato con altri abitanti del
villaggio saremmo stati avvisati molto prima di essere fermati. E così,
per il momento, mi sedetti accanto a Ina mentre col nastro adesivo
fissava una flebo salina (sigillata, senza ago, solo per fare scena)
nella piega del gomito di En. En era entusiasta della parte e cominciò a
tossire per fare le prove, una tosse polmonare che indusse Ina a fare
una smorfia in modo altrettanto teatrale: «Hai rubato le sigarette ai
chiodi di garofano di tuo fratello?»

En arrossì. L'aveva fatto per essere il più realistico possibile, disse.

«Ah, sì? Be', bada di non finire nella tomba prima del tempo, con la tua
pantomima.»

Nijon sbatté le porte posteriori, salì sul sedile del conducente, avviò
il motore e cominciammo il viaggio procedendo a scossoni verso Padang.
Ina disse a En di chiudere gli occhi. «Fai finta di dormire. Metti in
pratica il tuo talento teatrale.» Di lì a poco la sua respirazione prese
un ritmo regolare, con lievi sbuffi.

«È rimasto sveglio tutta la notte ad ascoltare musica» spiegò Ina.

«Mi stupisce che riesca comunque a dormire.»

«Uno dei vantaggi dell'infanzia. O della Prima Età, come la chiamano i
marziani: giusto?»

Annuii.

«Ne hanno quattro, ho capito bene? Quattro Età rispetto alle nostre
tre?»

Sì, era proprio come diceva lei. Di tutti gli usi e costumi delle Cinque
Repubbliche di Wun Ngo Wen, questo era l'aspetto che aveva affascinato
di più il pubblico terrestre.

Le culture umane riconoscono in genere due o tre fasi della vita:
infanzia e maturità; oppure infanzia, adolescenza, maturità. Alcuni
riservano alla vecchiaia una condizione particolare. L'usanza marziana,
invece, era unica e dipendeva dalla loro padronanza secolare di
biochimica e genetica. I marziani consideravano la vita umana divisa in
quattro parti, segnate da eventi mediati dalla biochimica. Il periodo
dalla nascita alla pubertà era l'infanzia. Dalla pubertà alla fine della
crescita fisica e all'inizio dell'equilibrio metabolico era
l'adolescenza. Dall'equilibrio al declino, alla morte o al cambiamento
radicale era la maturità.

E, oltre la maturità, l'età elettiva: la Quarta.

Secoli prima, i biochimici marziani avevano ideato un modo per
prolungare la vita umana di sessanta o settant'anni in media. La
scoperta, però, non era priva di svantaggi. Marte era un ecosistema
radicalmente limitato, dominato dalla scarsità di acqua e azoto. La
terra coltivata che era apparsa così familiare a Ibu Ina era un trionfo
di sottile e sofisticata bioingegneria. La riproduzione umana era stata
regolata per secoli, ancorata alle stime di sostenibilità. Altri
settant'anni aggiunti alla speranza di vita media rappresentavano una
crisi in fieri della popolazione.

E neppure il trattamento stesso per la longevità era semplice e
piacevole. Si trattava di una ricostruzione cellulare profonda. Veniva
iniettato nel corpo un cocktail di organismi virali e batterici
altamente modificati. I virus adattati eseguivano una sorta di
aggiornamento sistemico, aggiustando o correggendo le sequenze di DNA,
ricostruendo i telomeri, azzerando l'orologio genetico, mentre i
batteriofagi coltivati in laboratorio eliminavano metalli e placche
tossici e riparavano i danni fisici evidenti.

Il sistema immunitario resisteva. Il trattamento equivaleva, nella
migliore delle ipotesi, al decorso di sei settimane di una qualche
influenza debilitante: febbre, dolori articolari e muscolari, debolezza.
Certi organi cominciavano una specie di attività riproduttiva frenetica.
Le cellule della pelle morivano e venivano sostituite con una
successione furiosa; il tessuto nervoso si rigenerava spontaneamente e
in maniera rapida.

Il processo era debilitante, doloroso, con possibili effetti collaterali
negativi. La maggior parte dei soggetti riportava almeno la perdita
della memoria a lungo termine. Alcuni casi rari subivano una demenza
momentanea e un'amnesia irrecuperabile. Il cervello, ripristinato e
ricollegato, diventava un organo impercettibilmente diverso. E il suo
proprietario un essere umano impercettibilmente diverso.

«Hanno dominato la morte.»

«Non proprio.»

«Con tutta la loro saggezza, viene da pensare che avrebbero potuto
renderla un'esperienza meno spiacevole.»

Avrebbero sicuramente potuto alleviare il disagio superficiale della
transizione alla Quarta. Ma avevano scelto di non farlo. La cultura
marziana aveva incluso la Quarta Età nei suoi usi e costumi, con il
dolore e tutto il resto: il dolore era una delle condizioni restrittive,
un disagio tutelare. Non tutti sceglievano di diventare un Quarto. Non
solo la transizione era difficile, ma nelle loro leggi sulla longevità
erano state codificate rigide pene sociali. Qualunque cittadino marziano
era autorizzato a sottoporsi al trattamento, gratuitamente e senza
pregiudizi. Ma ai Quarti era proibito riprodursi; la riproduzione era un
privilegio riservato agli adulti. (Da duecento anni, nel cocktail della
longevità venivano aggiunti farmaci che provocavano la sterilizzazione
irreversibile per entrambi i sessi.) Ai Quarti non era permesso votare
alle elezioni del consiglio - nessuno voleva un pianeta governato da
anziani venerabili che avrebbero badato ai loro interessi. Tutte le
Cinque Repubbliche, però, avevano una specie di organo di revisione
giudiziaria, l'equivalente di una Corte Suprema, eletto
\emph{esclusivamente} dai Quarti. I Quarti avevano qualcosa in più e
altre cose in meno rispetto agli adulti, come gli adulti se paragonati
ai bambini. Più forti, meno giocosi; più liberi, ma anche meno liberi.

Ma non riuscivo a decifrare, né per Ina né per me stesso, tutti i codici
e i totem in cui i marziani avevano raccolto la loro tecnologia medica.
Gli antropologi ci avevano provato per anni, lavorando alla
documentazione archivistica di Wun Ngo Wen. Finché la ricerca non era
stata proibita.

«Ma ora abbiamo la stessa tecnologia» disse Ina.

«Solo alcuni di noi ce l'hanno. Spero che alla fine l'avremo tutti.»

«Mi chiedo se la useremo con altrettanta saggezza.»

«È possibile. I marziani l'hanno fatto, e i marziani sono umani quanto
noi.»

«Lo so. È possibile, certo. Ma tu cosa pensi, Tyler: ci
\emph{riusciremo}?»

Guardai En. Dormiva ancora. Forse stava sognando, gli occhi sotto le
palpebre chiuse guizzavano come pesci sott'acqua. Mentre respirava, le
narici si dilatavano e il movimento dell'ambulanza lo scuoteva da una
parte all'altra.

«Non su questo pianeta.»

ooo

A una quindicina di chilometri da Bukik Tinggi, Nijon bussò forte sul
divisorio tra noi e il posto del conducente. Quello era il nostro
segnale prestabilito: più avanti c'era un posto di blocco. L'ambulanza
rallentò. Ina si alzò in fretta e furia per prepararsi. Fissò una
maschera di ossigeno giallo fluorescente sul viso di En - il quale,
ormai sveglio, sembrava riconsiderare i meriti dell'avventura - e si
coprì la bocca con una mascherina di carta. «Fai presto» mi sussurrò.

E allora, contorcendomi, mi infilai nell'armadietto delle attrezzature.
Il coperchio batté sulle zeppe che facevano fluire un po' d'aria
all'interno, tra me e l'asfissia restava solo mezzo centimetro.

L'ambulanza si fermò prima che fossi pronto e urtai violentemente la
testa contro l'estremità stretta dell'armadietto.

«E adesso stai calmo» disse Ina - non capii se a me o a En.

Aspettai al buio.

I minuti passavano. Si sentiva un rimbombo lontano di gente che parlava,
impossibile da decifrare anche se avessi compreso la lingua. Due voci.
Nijon e uno sconosciuto. Una voce sottile, lagnosa, stridula. La voce di
un poliziotto.

«Hanno dominato la morte» aveva detto Ina.

``No'' pensai.

L'armadietto si stava riscaldando in fretta. Il sudore mi scivolava
sulla faccia, mi inzuppava la camicia, mi pizzicava gli occhi. Riuscivo
a sentire il mio stesso respiro. Immaginai che tutto il mondo potesse
sentirlo.

Nijon rispondeva al poliziotto sottovoce, deferente. Il poliziotto gli
urlava contro altre domande.

«Stai fermo ora, sta' fermo» sussurrò Ina con insistenza. En aveva
cominciato a far rimbalzare i piedi contro il materasso sottile della
barella, un'abitudine nervosa. Troppa energia per essere affetto da SDC.
Dal mezzo centimetro di luce sopra la testa, vidi divaricarsi la punta
delle dita di Ina, quattro ombre noccolute.

In quel momento le porte posteriori dell'ambulanza si aprirono
sbatacchiando e sentii l'odore di gas di scarico e di vegetazione
imputridita in un mezzogiorno estivo. Se allungavo il collo - piano
piano - riuscivo a vedere una sottile lama di luce all'esterno e due
ombre che potevano essere Nijon e un poliziotto o forse solo alberi e
nuvole.

Il poliziotto pretendeva qualcosa da Ina. La sua voce gutturale era
monotona, annoiata e minacciosa e mi faceva rabbia. Pensai a Ina ed En,
che si intimorivano o fingevano di intimorirsi davanti a quell'uomo
armato e a ciò che rappresentava. Lo facevano per me. Ibu Ina rispose
nella sua lingua in modo duro ma non provocatorio: «SDC, \emph{parole in
	minang}\ldots{} SDC.» Esercitava la sua autorità medica, testava la
suscettibilità del poliziotto, ricambiava paura con paura.

La risposta del poliziotto fu secca, chiedeva di perquisire l'ambulanza
o controllare i suoi documenti. Ina disse qualcosa di più energico o
disperato. Di nuovo, la parola SDC.

Volevo proteggere me stesso, ma ancora di più volevo proteggere Ina ed
En. Mi sarei arreso, piuttosto che vederli in pericolo. Arrendersi o
lottare. Lottare o fuggire. Rinunciare, se necessario, a tutti gli anni
da vivere in più che il farmaco marziano aveva donato al mio corpo.
Forse era quello il coraggio del Quarto, il coraggio particolare di cui
aveva parlato Wun Ngo Wen.

«Hanno conquistato la morte.» Ma no: come specie, terrestre o marziana,
in tutti i nostri anni su entrambi i pianeti avevamo solo messo in atto
dei rinvii di condanna. Nessuna soluzione finale.

Rumore di passi, stivali sul metallo. Il poliziotto stava salendo
sull'ambulanza. Potevo dire che fosse entrato dal modo in cui il veicolo
si abbassava sotto i colpi, ondeggiando come una nave cullata da un
lieve moto ondoso. Mi aggrappai al coperchio dell'armadietto. Ina si
alzò, disse di no con voce stridula.

Feci un respiro e mi preparai a saltar fuori.

Ma dalla strada provenne un altro rumore. Un veicolo ci superò rombando.
A giudicare dall'effetto doppler del gemito del motore sotto sforzo,
viaggiava ad alta velocità - una velocità cospicua, scioccante, da ``me
ne frego delle regole.''

L'agente ringhiò indignato. Il pavimento rimbalzò di nuovo.

Stropiccio di piedi, silenzio per il tempo di un battito, una portiera
sbattuta e poi il rumore dell'auto del poliziotto (immaginavo) mandata
su di giri e pronta a vendicarsi, con la ghiaia che si strappava via
dagli pneumatici in una grandine di rabbia.

Ina sollevò il coperchio del mio sarcofago.

Mi sedetti nel tanfo del mio stesso sudore. «Cos'è successo?»

«Era Aji, che arrivava dal villaggio. È un mio cugino. Correva al posto
di blocco per distrarre la polizia.» Era pallida ma sollevata. «Temo
stia guidando da ubriaco.»

«L'ha fatto per toglierci dalle balle la polizia?»

«Che espressione colorita. Sì. Siamo un convoglio, ricordi? Altre auto,
in costante contatto\ldots{} deve avere saputo che siamo stati fermati.
Rischia una multa o un ammonimento, niente di più.»

Respirai l'aria, che era dolce e fresca. Guardai En che mi fece un
sorriso incerto.

«Ti prego, presentami Aji quando arriviamo a Padang» dissi. «Voglio
ringraziarlo per aver finto di essere ubriaco.»

Ina alzò gli occhi al cielo. «Purtroppo Aji non fingeva. \emph{Era}
ubriaco. Un'offesa agli occhi del Profeta.»

Nijon ci guardò, fece l'occhiolino, chiuse le porte posteriori.

«Be', è stato spaventoso» disse Ina. Mi mise la mano sul braccio.

Mi scusai per averle lasciato correre quel rischio.

«Sciocchezze, adesso siamo amici. E non è così rischioso come pensi. Può
essere difficile trattare con la polizia, ma almeno sono uomini del
posto e legati a determinate regole, non come gli uomini di Jakarta, la
nuova Reformasi o qualunque sia il nome con cui si definiscono gli
uomini che hanno bruciato la mia clinica. E mi auguro che anche tu
rischieresti per noi, se fosse necessario. Lo faresti, Pak Tyler?»

«Certo che sì.»

Le tremava la mano. Mi guardò negli occhi. «Oddio, credo che tu dica sul
serio.»

No, non avevamo mai dominato la morte, solo messo in atto dei rinvii di
condanna (le pillole, le polverine, l'angioplastica, la Quarta Età) -
mettendo in scena la nostra convinzione che vivere più a lungo, anche
solo \emph{un po}' di più, avrebbe potuto ancora darci il piacere o la
saggezza che volevamo o che ci erano mancati. Nessuno torna a casa dopo
un triplo bypass o una cura per la longevità aspettandosi di vivere per
sempre. Persino Lazzaro lasciò la tomba sapendo che sarebbe morto per la
seconda volta.

Ma lui si alzò e camminò. Camminò con gratitudine. Ero grato anch'io.